\documentclass{article}
\usepackage{geometry}
\usepackage{amsmath}
\usepackage{amsfonts}
\usepackage{graphicx}
\usepackage{caption}
\usepackage{subcaption}
\usepackage{listings}
\usepackage{color}
\usepackage{hyperref}
\usepackage{tikz}
\usepackage{setspace}
\usepackage{booktabs}
\usepackage{siunitx}
\usetikzlibrary{3d}
\usepackage{pgfplots}
\pgfplotsset{compat=1.18}
\geometry{a4paper, margin=1in}
\doublespacing
\begin{document}

\begin{titlepage}
    \centering
    \vspace*{2cm} % Add vertical space at the top
    
    {\Huge \textbf{A Survey of Reinforcement Learning in the Nuclear Industry}} \\
    
    \vspace{2cm}
    
    {\Large \textbf{Danen M. Pruitt}} \\
    
    \vspace{1.5cm}
    
    ENS, AFIT \\
    OPER 646: Reinforcement Learning \\
    May 2026
    
    \vfill % Pushes the following text to the bottom
    
\end{titlepage}

\newpage
\tableofcontents
\newpage

\section{Introduction}
Reinforcement learning (RL) has undergone a significant shift by transitioning from a niche machine learning process into a foundational technology for autonomous decision-making in complex, real-world systems. Unlike supervised learning, which relies on static, labeled datasets to predict patterns, RL operates on an active learning principle. By interacting with a dynamic environment, an RL agent learns through a process of trial and error utilizing a feedback loop where actions are evaluated against a reward function. This allows the agent to optimize its behavior and develop strategies that maximize cumulative rewards over time. 
\newline
This environmental interaction is precisely what makes RL so powerful in unpredictable, real-world settings. Traditional control systems often rely on a single, static policy that may become obsolete as environmental conditions shift. However, an RL agent is designed to continually adapt its policy to the ever-changing state of its surroundings. This capability is crucial for bridging the gap between simulation and reality, which is the notoriously difficult task of transferring an agent from a sterile simulation to the friction and noise of reality. Furthermore, the recent integration of RL with Large Language Models (LLMs) has provided these algorithms with a significant boost by combining high-level logical reasoning and long-horizon planning with low-level motor or logic control. Thus, RL has evolved from creating simple behavioral agents into developing sophisticated, effective collaborators.
\newline
As these technological barriers fall, the deployment of RL has accelerated across diverse sectors. In healthcare, it is being utilized to pioneer Personalized Medicine, where RL agents design dynamic treatment regimes that adjust in real-time to a patient’s physiological response. Beyond clinical settings, the application of RL is improving soft robotics, autonomous vehicle navigation, and the management of smart electrical grids. It has even proven essential in industrial sustainability, where it is used to optimize HVAC controls and energy consumption in large-scale facilities by balancing comfort with extreme operational efficiency.
\newline
The primary objective of this survey is to examine how these RL applications are and could be applied and purpose-built to the specific complexities of the nuclear industry. The scope ranges from initial core loading and fuel assembly optimization to the complex logistics of waste management and chemical reprocessing. To ensure a high standard of technical relevance, the papers analyzed in this review were selected from Q1 and Q2 journals published between 2020 and the present. These selections represent the current state-of-the-art methods for applying reinforcement learning to the unique safety, precision, and regulatory requirements of the nuclear industry.

\newpage
\section{Methodology}
These 12 papers were selected specifically from journals that are ranked Q1 or Q2 in the Scimago Journal Ranker or was a top 20 journal or conference proceeding as ranked by Google Scholar. The most efficient process used to find these articles was to search Google Scholar using the following search string format:
\begin{align}
    \text{source:"Journal or Conference", "Reinforcement Learning and other key words"}
\end{align}
This allowed for certainty that the articles and sources found were from reputable sources while searching for specific use cases. The journals and conferences searched from included:
\begin{itemize}
    \item Nuclear Engineering and Design (Q1)
    \item Nuclear Engineering and Technology (Q2)
    \item Annals of Nuclear Energy (Q1)
    \item Progress in Nuclear Energy (Q1)
    \item Applied Energy (Top 20 - Engineering)
    \item Energy (Q1)
    \item IEEE Transactions on Nuclear Science (Q1)
    \item IEEE International Conference on Robotics and Automation (Top 20)
    \item American Nuclear Society Winter/Annual Meetings (Top 20)
\end{itemize}
Because there was not enough research on any one portion of the nuclear energy industry, the following groups formed:
\begin{itemize}
    \item Nuclear Core and Fuel Assembly Design Optimization 
    \item Autonomous Reactor Control and Operations 
    \item Spent Fuel Management and Broader Plant Management
\end{itemize}
The goal of looking at this rather large application domain was to see how the basics as well as the specifics of RL are used by various professionals in the nuclear industry. The primary use was found in design and fuel loading which is a very difficult task to plan for and optimize which will be discussed later. 

\newpage
\section{Summaries}
\subsection{Nuclear Core and Fuel Assembly Design Optimization}
\textbf{Large-scale Design Optimization of Boiling Water Reactor Bundles with Neuroevolution}
\newline
The design and optimization of Boiling Water Reactor (BWR) fuel bundles is a massive, high-dimensional combinatorial challenge characterized by heavy computational costs and strict safety constraints. To overcome the limitations of traditional stochastic or purely neural network-based algorithms, this study introduces a novel neuroevolution approach that combines the strengths of Deep Reinforcement Learning (RL) and Evolution Strategies (ES). Specifically, the researchers utilize the Proximal Policy Optimization (PPO) RL algorithm to narrow the vast design space into feasible regions by learning optimal placements for different enrichment assemblies. Once these regions are identified, Evolution Strategies are employed to thoroughly explore the surrounding areas, preventing the optimizer from getting trapped in local optima. Unlike previous studies that simplified the problem to 2D models, this methodology optimizes the full 3D 10x10 bundle using the high-fidelity CASMO4 and SIMULATE3 industry codes.
\newline
To tackle the immense computational burden of evaluating over 40 distinct design constraints—including power peaking factors (PPF), cold shutdown margins, and control rod insertions—the researchers implemented a highly efficient "divide-and-conquer" reward shaping strategy. This approach strategically groups constraints so that the algorithm first satisfies simple, fast-to-calculate geometric rules before progressing to computationally expensive simulations. The results demonstrated that this RL-informed ES neuroevolution hybrid was immensely successful and highly computationally efficient, outperforming standalone RL by discovering approximately 100 unique, feasible 3D bundle patterns compared to the mere 14 found by RL alone. The best-optimized bundle generated by this framework achieved a highly efficient average burnup while successfully minimizing radial power peaking factors and strictly adhering to all complex operational limits.
\newline
\textbf{Assessment of Reinforcement Learning Algorithms for Nuclear Power Plant Fuel Optimization}
\newline
This paper also discusses the combinatorial and multi-objective problem posed by fuel core reloads. Traditionally this problem is tackled using hand-designed heuristics or stochastic methods like genetic algorithms and simulated annealing. The goal was to improve by introducing an approach using deep reinforcement learning, specifically the proximal policy optimization (PPO) algorithm, to autonomously design optimal core configurations. In this RL framework, the agent sequentially constructs loading patters by selecting and placing fresh and burning fuel assemblies, receiving a reward based on an objective function that minimized the LCOE while still strictly enforcing the critical safety limits.
\newline
The primary focus of this study was extensive analysis of the governing hyper-parameters of the PPO algorithm to understand the exploration-vs-exploitation trade off. The researchers discovered that deep RL achieves the highest sample efficiency and stability when applied similarly to a Gaussian process, specifically by restricting the number of patterns generated per episode to one and carefully reducing step counts and entropy coefficients until learning is maximized. Ultimately this rigorously tuned framework demonstrated strong generalizability on mathematical benchmarks as well as created optimized core designs that offer high economic benefit. 
\newline
\textbf{Physics-informed Reinforcement Learning Optimization of PWR Core Loading Pattern}
\newline
This paper sheds more light on the complexity of a core reload by explaining that the problem is a highly complex combinatorial problem with a massive search space of up to $10^{38}$ possible configurations, thus not explicitly solvable. The goal was to minimize the levelized cost of electricity (LCOE) while strictly adhering to the strict safety and operational constraints in a pressurized water reactor. To accomplish this the researchers developed a physics informed multi-objective deep reinforcement learning approach called PEARL. Using various methods, the problem was strategically divided into a two-step curriculum learning process. The agent was first trained to satisfy highly correlated constraints such as boron concentration and peak pin power, and then it concurrently balanced the most difficult trade-offs such as LCOE, cycle length, and rod integrated peaking factor.
\newline
After training the limits of this framework were tested by applying it to restricted design scenarios that aggressively push beyond standard expert designs. While standard single-objective RL methods proved completely useless in finding feasible and safe solutions to these highly complex situations, the physics informed two step approach of PEARL successfully discovered a large number of diverse loading patterns. This demonstrated substantial economic potential, generating optimized core patterns that offer estimated cost savings between two and three million per year per plant when compared to the baseline designs.
\newline
\textbf{Multi-Objective Reinforcement Learning-based Approach for Pressurized Water Reactor Optimization}
\newline
This paper adds additional information to the complexity of nuclear core reloading patterns. However, it separates this implementation from other methods by exclusively describing them as multi-objective methods. It mentions the PEARL (pareto envelope augmented with reinforcement learning) algorithm. In this paper the researchers introduced a constrained variant called C-PEARL that utilizes curriculum learning where the agent first learns to consistently generate feasible, safe solutions before progressing to the more difficult task of optimizing the multi-objective trade-offs within the feasible space.
\newline 
To evaluate this and other variants of PEARL, such as PEARL-NdS, they were compared to classical mathematical benchmarks and practical PWR core optimization scenarios featuring two and three competing objectives. The RL agent’s specific performance was systematically compared to multi-objective stochastic optimization methods such as the Non-dominated Sorting Genetic Algorithm (NSGA-III), Simulated Annealing (SA), and Tabu Search (TS) which are heavily utilized in the industry.  Results of these comparisons showed that the PEARL-NdS variant combined with curriculum learning systematically outperformed the classical industry methods across multiple metrics including the sheer cardinality of optimal solutions found. Ultimately this approach successfully captured the topological structure of the optimization problem, offering core designers a highly efficient tool to automatically uncover a superior range of optimal fuel loading trade-offs.
\newline 
\newline\textbf{Impact of including fuel performance as part of Core Reload Optimization: application to power uprates}
\newline
This paper investigates the integration of critical fuel performance metrics directly into the reinforcement learning-based optimization of core loading patterns for pressurized water reactors undergoing power uprates. Reinforcement Learning algorithms along with deep learning ensembles were applied because traditional physics simulation codes are extremely computationally expensive for the tens of thousands of evaluations required during the optimization process.  The models were employed as fast-running surrogates which along with the shorter computational time provide robust uncertainty estimates. Testing across various pressurized water reactor and uprate scenarios, the study found that utilizing a low fidelity version of the surrogate models offered the most practical trade-off between computational speed and fuel performance gains. Ultimately, while integrating these models into the multi-objective reward function only yielded marginal improvements to the fuel performance, it led to a significant enhancement in the safety limits for hydraulics and neutronics. This demonstrated that reinforcement learning agents can effectively exploit intrinsic physical correlations between core constraints when coupled with the core reload design process. 
\newline
\textbf{Physics-informed reinforcement learning optimization of nuclear assembly design}
\newline
This paper explores the application of deep reinforcement learning (RL) algorithms—specifically Deep Q-Learning (DQN) and Proximal Policy Optimization (PPO)—to the problem of optimizing Boiling Water Reactor (BWR) fuel assemblies. Because arranging nuclear fuel and burnable poison rods to maximize fuel cycle length while strictly adhering to safety constraints (such as power peaking factors) yields a massive computational search space, the authors introduce a "physics-informed" optimization approach. By embedding nuclear engineering heuristics and expert knowledge directly into the RL reward system as "game rules," the trained AI agent learns to navigate the design space efficiently without wasting time evaluating unsafe designs. When tested on both a low-dimensional 6x6 assembly and a high-dimensional 10x10 assembly using the CASMO-4 physical simulator, the RL methodology significantly outperformed traditional stochastic optimization methods like genetic algorithms and simulated annealing. Ultimately, the RL framework successfully discovered four to five times more feasible assembly patterns and demonstrated outstanding computational speed and efficiency, proving to be a highly robust decision-support tool for advancing nuclear fuel design


\subsection{Autonomous Reactor Control and Operations}
\textbf{Small modular reactor reinforcement learning framework: Automating reactor core startup}
\newline
A crucial technology in moving toward carbon neutrality is small modular reactors (SMRs) but they require advanced automation in order to reduce human error and operator workload. This paper discusses a particularly difficult manual operation in the reactor core start up which requires delicate sequential manipulation of control rods to safely transition the core from a subcritical state to a critical state without causing power runaway. Researchers addressed this by developing a reinforcement learning framework capable of autonomously managing the SMR core startup process. The agent was trained in a S-SMRSim environment which is a dynamic simulation model validated against existing experimental data. The simulator allowed the agent to continuously interact with the environment without the massive computational overhead of traditional physics codes.
\newline
Within the simulated environment the researchers trained a deep Q-network (DQN) algorithm and a Proximal Policy Optimization (PPO) algorithm to control the reactor. Sequentially they were tasked with deciding to withdraw, insert, or halt control rods, relying on state information comprised of core power and rate of power change. To better guide learning, a two part penalty based reward system was implemented, penalizing the agent for deviating from the target critical power level and for exhibiting unsafe power escalations. The results showed that both algorithms successfully automate the core start up reliably increasing power. In the end the DQN algorithm exhibited faster convergence due to its update mechanisms being better suited for the discrete decisions in the step by step manipulation of rods.
\newline
\textbf{Improvement of a reinforcement learning model to assist operators in controlling a nuclear power plant under abnormal operating conditions}
\newline
This paper addresses critical operational challenges in advanced nuclear power plants (NPPs) by refining an AI-based Counter-Measure Proposal Module (CMPM). The CMPM utilizes reinforcement learning to assist plant operators during anomalous events by autonomously suggesting rapid, coordinated countermeasures aimed at restoring main plant parameters—specifically reactor power and core-outlet temperature—back to normal steady-state levels. While a previous iteration of the CMPM successfully mitigated anomalies, it exhibited two significant execution flaws that could confuse human operators: it occasionally recommended unnecessary corrective actions even when the plant was already operating normally, and it sometimes proposed conflicting operations for parallel components that serve the exact same function, such as signaling one auxiliary water pump to increase its rotation rate while simultaneously commanding the other to decrease.
\newline
To resolve these practical issues, the researchers structurally improved the reinforcement learning model by meticulously redesigning its reward functions. They introduced new mathematical penalty terms into the training algorithm—specifically using Proximal Policy Optimization (PPO)—to heavily discourage the agent from taking action when plant deviations are negligible, and to explicitly penalize the model for proposing counterproductive, opposing actions for related auxiliary pumps. By training the agent on computationally efficient surrogate models representing the High Temperature engineering Test Reactor, the authors demonstrated that the newly formulated reward function successfully eliminated the illogical and unnecessary operational suggestions. Ultimately, the improved CMPM maintained its high effectiveness in stabilizing the plant during disturbances, providing a safer, leaner, and more reliable decision-support tool for operators without compromising its mitigating capabilities.
\newline
\textbf{Reinforcement learning: Application to PWR's core power level control}
\newline
This paper presents the design and implementation of an Adaptive Dynamic Programming (ADP) based controller to regulate the core power level of Pressurized Water Reactors (PWRs). Nuclear power plants are inherently complex, nonlinear, and time-varying systems where operational characteristics fluctuate significantly due to factors like fuel burnup, temperature feedback, and equipment degradation. Because of these continuous changes, traditional control strategies with fixed gains often become ineffective or unstable over the plant's lifecycle, especially during demanding load-following operations. To address this, the researchers introduce a model-free ADP control strategy that does not require a pre-existing dynamical model of the reactor. Instead, the control algorithm transforms the reactor's unknown nonlinear dynamics into a normal form and uses real-time state measurements to iteratively learn and approximate an optimal control law, enabling it to dynamically adapt to the system without needing to reconstruct unmeasurable states.
\newline
To evaluate the efficacy of the reinforcement learning framework, numerical simulations were conducted subjecting the controller to various load-following patterns, such as power descents to 50\% capacity and power lifts back to full power. The ADP controller was benchmarked against traditional Proportional-Integral-Derivative (PID) and Linear Quadratic Gaussian (LQG) controllers. The results demonstrated that the ADP method significantly outperformed the conventional controllers; it smoothly tracked the power demand references without the noticeable undershoots, overshoots, or steady-state errors exhibited by the PID and LQG models. Furthermore, the proposed controller required less control effort—meaning smaller and fewer control rod adjustments—and proved highly robust against both matched and unmatched system uncertainties and external disturbances. The study confirms that ADP-based reinforcement learning provides a stable, adaptive, and highly efficient solution for managing the complex, time-varying dynamics of nuclear reactors.


\subsection{Spent Fuel Management and Broader Plant Management}
\textbf{Applications of deep reinforcement learning in nuclear energy: A review}
\newline
This paper provides a comprehensive survey of how Deep Reinforcement Learning (DRL) is being adapted for use within the nuclear energy sector. The authors begin by outlining the theoretical foundations of DRL, which combines the high-dimensional perceptual capabilities of deep learning with the dynamic decision-making and optimal strategy-learning strengths of reinforcement learning. Before examining nuclear-specific use cases, the review analyzes successful DRL implementations across other complex industrial fields—such as autonomous driving, robotics, wind and solar energy management, and power grid dispatching—to establish a baseline for its capabilities in task planning, resource scheduling, and system control.
\newline
Within the nuclear industry, the review categorizes current DRL applications into several key areas: autonomous reactor control (such as startup and shutdown operations), operation and maintenance strategy optimization, physical design optimization (like nuclear fuel assembly configurations), and safety and fault diagnosis. The authors emphasize that while DRL offers robust adaptability and the ability to find optimal solutions without requiring extensive mathematical modeling beforehand, its deployment in nuclear environments faces rigid safety and reliability constraints. Key challenges that future researchers must overcome include the inefficiency of data sampling, difficulties in designing accurate reward functions, and the "simulation-to-reality" gap when transferring algorithms from virtual plant simulators to actual physical systems. Despite these hurdles, the paper concludes that DRL holds significant potential to advance the automation, safety, and economic efficiency of next-generation nuclear energy systems.
\newline
\textbf{Potential applications of artificial intelligence in spent nuclear fuel reprocessing research: A review}
\newline
This review paper examines the transformative role of artificial intelligence (AI)—including machine learning and deep learning—in overcoming the significant challenges associated with spent nuclear fuel reprocessing. Traditional reprocessing research is heavily constrained by extreme radiation hazards, a scarcity of expensive experimental data, and the intricate, nonlinear mechanisms of multiphase chemical separations. To address these bottlenecks, the authors explore AI's emerging applications across several critical operational domains. AI is being utilized for data-driven chemical reaction modeling to efficiently uncover hidden kinetic pathways, and for the intelligent design of new materials, such as screening highly selective extractants and robust, radiation-resistant adsorbents. Additionally, AI enhances autonomous process control and real-time monitoring—using techniques like deep reinforcement learning and machine learning-assisted spectroscopy—while also improving data analysis for precise spent fuel characterization and rapid quality control.
\newline
Despite these promising advancements, the paper emphasizes that widespread AI deployment in the nuclear sector faces substantial hurdles regarding data, interpretability, and safety. High-quality experimental data is difficult to obtain, often leading to an overreliance on simulated data that may not accurately capture real-world complexities or rare failure conditions. Furthermore, the "black box" nature of deep learning models poses a major trust barrier in a safety-critical industry, necessitating the development of Explainable AI (XAI) to ensure that the AI's decision-making logic is transparent and verifiable by human experts. To guarantee reliability and defend against cybersecurity threats, the authors advocate for hybrid methodologies that deeply integrate physics-based first principles with data-driven AI, supported by digital twin technologies and rigorous "human-in-the-loop" oversight. The review concludes that developing these robust and interpretable AI systems holds the potential to significantly enhance the safety, automation, and economic viability of the closed nuclear fuel cycle.
\newline
\textbf{Reinforcement learning for vulnerability assessment of physical protection systems}
\newline
This paper addresses the growing computational challenges of identifying security vulnerabilities in Physical Protection Systems (PPS) that safeguard critical infrastructure. Traditional vulnerability assessment tools, such as the classic Adversary Sequence Diagram (ASD) or basic heuristic searches, often struggle with high computational complexity and analytical burden when attempting to map the vast space of potential adversarial paths in realistic 3D environments. To overcome this, the authors propose a novel Q-learning reinforcement learning framework built upon an extended ASD model. By mapping interconnected security zones and protective elements (like doors, fences, and cameras) into a defined state space, the researchers trained "adversary" and "response force" agents to simulate physical intrusions. The adversary agent is trained to find the most exploitable attack routes by maximizing rewards for reaching the target, while receiving penalties for lingering in monitored areas or sustaining combat injuries. Conversely, the response force agent is trained to find the optimal defensive routes to intercept and neutralize the threat as quickly as possible.
\newline
To validate the methodology, the researchers simulated a 3D hypothetical laboratory facility equipped with a robust layout of detection and delay elements, such as access control gates and video surveillance. After approximately 54 minutes of upfront offline training, the Q-learning algorithm successfully converged and identified the facility's most vulnerable adversarial intrusion paths. When benchmarked against a traditional heuristic search algorithm, the proposed Q-learning method proved vastly superior in real-time efficiency: it required only 57 node traversals and 275 milliseconds to identify the vulnerable path, compared to the algorithm's 145,521 node traversals and roughly 752 seconds. The study demonstrates that integrating reinforcement learning with extended ASD models can significantly reduce the computational cost of pathfinding in large-scale 3D environments, providing a fast and highly effective tool for optimizing security strategies and response deployments.

\section{Comparative Analysis and Discussion}
Across the surveyed papers, RL approaches were carefully tailored to match the specific action/state spaces, problem sizes, and operational requirements of varying nuclear applications. Algorithms that estimate the value of states and/or state-action pairs were heavily utilized for discrete and low dimensional tasks. Q- learning for example was applied in the physical protection systems vulnerability assessments while the small modular reactor core startup empoyed deep q-networks since there was a small discreate action space but larger state space which proved highly effective. 
\newline
Other algorithms applied policy based or actor critic models where high-dimensional or continuous state-action spaces were present. Proximal policy optimization is the dominant algorithm for the large scale combinatorial optimization of fuel loading patterns and bundle design. This in addition to the physics-informed reward shaping and curriculum learning strategies were highly effective in handling the astronomically large search spaces and complex multi-objective trade-offs of these problems. Additionally the integration of a neuroevolution framework, narrowing the vast search space achieved robust optimization for the bundle problem.
\newline
Other algorithms included adaptive dynamic programming for the core power level control and PEARL to balance competing objectives in core reload and facility optimization. Overall the slight additions and enhancements to various proven RL algorithms has proven very successful in handling the unique challenges of the nuclear industry with the most common problem in this review being the loading of the fuel assemblies. Below is a table summarizing some of the applications. 
\begin{table}[htbp]
    \centering
    \caption{Comparison of Reinforcement Learning (RL) Applications in Nuclear Engineering}
    \label{tab:rl_applications}
    \resizebox{\textwidth}{!}{%
    \begin{tabular}{lllll}
        \toprule
        \textbf{Application Domain} & \textbf{Algorithm Used} & \textbf{State / Action Space} & \textbf{Problem Size \& Scope} & \textbf{Effectiveness \& Key Features} \\ \midrule
        \begin{tabular}[c]{@{}l@{}}BWR/PWR Core\\ Loading Pattern\end{tabular} & \begin{tabular}[c]{@{}l@{}}PPO, PEARL,\\ Neuroevolution\end{tabular} & \begin{tabular}[c]{@{}l@{}}High-dimensional /\\ Discrete (Combinatorial)\end{tabular} & \begin{tabular}[c]{@{}l@{}}Massive ($10^{31}$ combinations\\ for $10\times10$ arrays)\end{tabular} & \begin{tabular}[c]{@{}l@{}}High. PPO finds 4--5$\times$ more feasible\\ patterns than SA/GA. \end{tabular} \\ \addlinespace
        SMR Core Startup & DQN, PPO & \begin{tabular}[c]{@{}l@{}}Continuous states (power rate) /\\ Discrete actions (rod motion)\end{tabular} & \begin{tabular}[c]{@{}l@{}}Low-to-Medium ($10^{-8}$ to\\ $10^{-2}$ power transition)\end{tabular} & \begin{tabular}[c]{@{}l@{}}High. DQN specifically excels at targeted\\ learning over PPO due to frequent\\ per-timestep parameter updates.\end{tabular} \\ \addlinespace
        \begin{tabular}[c]{@{}l@{}}Reactor Power Control\\ (Load Following)\end{tabular} & \begin{tabular}[c]{@{}l@{}}ADP (Adaptive Dynamic\\ Programming)\end{tabular} & Continuous / Continuous & \begin{tabular}[c]{@{}l@{}}Complex, unknown nonlinear\\ dynamics\end{tabular} & \begin{tabular}[c]{@{}l@{}}High. Model-free approach that relies\\ only on real-time state measurements.\end{tabular} \\ \addlinespace
        \begin{tabular}[c]{@{}l@{}}PPS Vulnerability\\ Assessment\end{tabular} & Q-learning & \begin{tabular}[c]{@{}l@{}}Discretized 3D blocks /\\ Discrete movements\end{tabular} & \begin{tabular}[c]{@{}l@{}}Medium (3D hypothetical\\ facility mapping)\end{tabular} & \begin{tabular}[c]{@{}l@{}}High. Reduces path search time to\\ 275ms after $\sim$54 minutes of offline training.\end{tabular} \\ \addlinespace
        \begin{tabular}[c]{@{}l@{}}Abnormal Condition\\ Countermeasures\end{tabular} & PPO & \begin{tabular}[c]{@{}l@{}}Continuous plant parameters /\\ Discrete component operations\end{tabular} & \begin{tabular}[c]{@{}l@{}}Medium (HTTR anomaly\\ recovery)\end{tabular} & \begin{tabular}[c]{@{}l@{}}High. Successful stabilization. Required\\ strict reward shaping to prevent counter-\\ productive or unnecessary actions.\end{tabular} \\ \bottomrule
    \end{tabular}%
    }
\end{table} 
\newline
Some strengths of the RL approaches are the ability to navigate large search spaces and escaping local optima in the complex combinatorial problems by shaping rewards according to the physics of the problem and using curriculum learning. The agents also act as near instant decision support tools once trained, which is crucial for the complicated control problems such at core startup and power control. However, the training of these models can be computationally expensive especially since getting sufficient training data through high grade physics simulators is the primary method of data creation. The common issues associated with RL still remain such as hyperparameter sensitivity and reward hacking if the rewards are not properly defined.
\newline
The biggest gap so far is the lack of real world application, testing, and deployment of these algorithms. In all of the sources, the algorithms were only tested in simulated environments. This could be due to the lack of generalization/standardization of nuclear codes across the industry as well as the strict safety regulations that are present. However, researchers are doing their best to get as close as possible with surrogate models and highly validated simulations.
\newline
Looking forward, the application of RL in the nuclear industry is influencing the broader RL field by forcing algorithms to adapt to the Zero tolerance safety constraints and heavily constrained multi-objective environments. This is building the foundation for robust and adversarial RL algorithms that can be applied to other safety critical industries. The combination of surrogate models with these RL methods also drastically increases sample efficiency and would be a major area of research in the near future.
\section{Conclusion}
Overall, RL methods along with the various enhancements are being studied heavily in nuclear industry with strong emphasis in the larger problems that often need expert input. These agents allow for the efficient creation of optimized designs and control strategies/operations increasing the ability of experts to make informed decisions and explore possbile solutions with more confidence and detail. Overall these papers, both application and review, contribute greatly to the opportunity RL has to revolutionize the industry by providing powerful tools to navigate the complex, high-dimensional, and safety-critical challenges of nuclear energy systems. The continued development and refinement of these algorithms in this industry, along with efforts to bridge the gap between simulation and real-world deployment, will be crucial for realizing the full potential of the safety, efficiency, and sustainability of nuclear power generation.



\newpage
\section{Bibliography}
\begin{itemize}
    \item Radaideh, M. I., Forget, B., \& Shirvan, K. (2021). Large-scale design optimisation of boiling water reactor bundles with neuroevolution. \textit{Annals of Nuclear Energy}.
    \item Liu, Y., Wang, B., Tan, S., Li, T., Lv, W., Niu, Z., Li, J., Gao, P., \& Tian, R. (2024). Applications of deep reinforcement learning in nuclear energy: A review. \textit{Nuclear Engineering and Design}, 429, 113655.
    \item Yoshikawa, M., Seki, A., Okita, S., Takaya, S., \& Yan, X. (2025). Improvement of a reinforcement learning model to assist operators in controlling a nuclear power plant under abnormal operating conditions. \textit{Nuclear Engineering and Design}, 444, 114350.
    \item Seurin, P., \& Shirvan, K. (2024). Multi-Objective Reinforcement Learning-based Approach for Pressurized Water Reactor Optimization. Manuscript, March 19, 2024.
    \item Radaideh, M. I., Wolverton, I., Joseph, J., Forget, B., Roy, N., \& Shirvan, K. (2021). Physics-informed reinforcement learning optimization of nuclear assembly design. \textit{Nuclear Engineering and Design}, 372, 110966.
    \item Li, T., Chen, Q., Zhou, J., Sun, Y., Jia, Y., Bai, Y., Guo, J., Zhang, Q., Wang, X., Yao, F., Cao, Z., Wang, W., Liu, F., Yan, T., \& Zheng, W. (2026). Potential applications of artificial intelligence in spent nuclear fuel reprocessing research: A review. \textit{Nuclear Engineering and Technology}, 58(4), 104060.
    \item Zhu, X., Guo, P., Zhang, Y., Ma, H., \& Yang, J. (2026). Reinforcement learning for vulnerability assessment of physical protection systems. \textit{Annals of Nuclear Energy}, 233, 112298.
    \item Boubacar Kirgni, H. (2025). Reinforcement learning: Application to PWR's core power level control. \textit{Progress in Nuclear Energy}, 180, 105644.
    \item Bae, S. J., Son, H. H., Lee, Y., \& Yang, J. (2025). Small modular reactor reinforcement learning framework: Automating reactor core startup. \textit{Nuclear Engineering and Technology}, 57(3), 103247.
    \item Seurin, P., \& Shirvan, K. (2024). Assessment of Reinforcement Learning Algorithms for Nuclear Power Plant Fuel Optimization. \textit{Applied Intelligence}.
    \item Seurin, P., \& Shirvan, K. (2024). Physics-informed Reinforcement Learning Optimization of PWR Core Loading Pattern. Manuscript, March 22, 2024.
    \item Seurin, P., Halimi, A., \& Shirvan, K. (2024). Impact of including Fuel Performance as part of Core Reload Optimization: application to power uprates. Manuscript, November 23, 2024.
\end{itemize}
\end{document}